\section{Introduction}

\iffalse

[ada-ppi construction]
- the importance of understanding PPI Networks as well as the complexes and pathways within.
- The expensiveness of experiments compared to computational methods, difference of PPI prediction vs functional module detection. a function module can be a protein complex or  pathway.
- problem definition + not a link prediction problem. biological information + structural information. over-smoothing and under smoothing

[2022-from-comunities]
- the importance of Proteins, PPI and their responsibilities in biological functions
- the experimental methods are expensive and critizied for high number of false positive and false negatives. (Has mentioned some of the experimental methods name)
- what are PPI networks, the complexes and clusters, go terms and the both structural(topology) and functional information. the go explainations can be useful.
- clusters and protein complexes (the relation). also the fact that some methods only consider the topology of network which is insufficent for detecting function modules.
- because of mistakes in PPI netowrks solely depending on network topology makes model to make incorrect predictions.
- talking about their method, what are the cons and the evaluating to SOTA methods. we are the best.

[2021-identifying]
- proteins are important, they act in complexes, researches now focus on complexes. the study of complexes is important.
- some of the important biological processes that done by protein complexes. telling about disseases and the PPIs. discovery of new pathways and proteins functionality.new methods for desiease control

[2021-TADW-SC]
- why proteins are important. they work together. what are the PPI networks. the same properties as complex networks such as small-world and high degree of clustering coefficient. many efforts.
- high cost and time requirements. challenging topic and one way is clustering.
- cluster definition. two main approaches --> topological and topology+information
- why using protein informations? what are the sources and GO explain.
- a small introduction to the method and the challenges
- the challenge of fusing topo + features. what are the best features 
- what are the each sections.


[2022-from local]
The pairwise physical interaction of proteins within the PPI data suggests a network representation where nodes are the proteins and links are the interactions among the proteins. Exploiting network structure with network analysis tools on such data has shifted the PPI analysis to the network level. Besides, the existence of protein complexes justifies the high-degree clusters within the PPI network [9]. PPI networks inherit both topological and functional information [1]. The first term refers to the physical interaction describing the arrangements of the nodes in the network, and is associated with the densely connected proteins namely communities. The latter explains the biological function of proteins that are achieved by groups of proteins that bind each other and shape protein complexes. The complexes are explained by the annotations available in Gene Ontology (GO) [11, 12] databases. GO provides a specific definition of protein functions and it is one of the main resources of biological information. GO provides a structured and controlled vocabulary of terms, which are subdivided into three non-overlapping ontologies: Molecular Function (MF), Biological Process (BP) and Cellular Component (CC) [13]. We utilize GO terms to enrich PPI networks with functional properties of proteins.


[2021-tadwsc]
In this approach, a protein complex is modelled as a set of clusters, aka communities, and each cluster is a complex of proteins that are interacting to perform the same biological function. Generally, in complex networks, clusters or communities are groups of nodes that contain the maximum number of intra-links and the minimum number of interlinks [3]. The task of community detection is defined as the identification of such node groups using the network’s properties [5]. In particular, the main goal of network clustering in the PPI network is to find the sub-networks of the original network, such that proteins in each cluster share similar topological and functional properties [3]. Therefore, we are able to figure out the function of a protein using the cluster it belongs to. In general, we can deal with the network clustering problems in PPI networks using two main approaches. One of these approaches is the use of protein network topological information. These techniques consider the PPI network’s topological information to explore and identify protein clusters. There are two main classes in this category: divisive and agglomerative. Divisive methods use a top-down technique and try to divide the original network into sub-networks [6]. On the other hand, agglomerative methods start with very small-sized proteins and work their ways by merging them until the final clusters are obtained [7].
\fi


Proteins play a fundamental role in biological processes and cellular functionality within organisms, mediating essential tasks such as DNA transcription, cell cycle regulation, and enzymatic activities. It is crucial to understand that proteins rarely act in isolation. instead, they accomplish these biological goals by interacting with one another. In other words, almost every biological process within an organism requires the coordinated interaction of multiple proteins. Consequently, in recent years, biologists have shifted their focus toward studying protein functions within collective groups rather than analyzing each protein individually.

These groups of interacting proteins, which work together to perform specific biological functions, are known as functional modules and are often referred to as protein complexes. Importantly, proteins are often multi-functional and may participate in multiple biological processes simultaneously, which frequently results in overlapping functional modules. The discovery of previously unknown protein complexes, or the association of uncharacterized proteins with known complexes, is vital for identifying new biological roles and annotating the function of unknown proteins. Such insights are critical for advancing function discovery and uncovering the complex molecular mechanisms underlying various diseases. 

Protein-protein interactions (PPIs) are typically modeled as networks, where nodes represent individual proteins and edges denote their physical interactions. These PPI networks exhibit topological characteristics common to real-world complex networks such as the small-world property and a high clustering coefficient. In recent years, high-throughput experimental techniques have been developed to construct these networks, providing a comprehensive understanding of protein interworkings. However, despite their utility, these experimental methods are often criticized for being costly and time-consuming, and for their inherent susceptibility to high rates of false positives and false negatives. Given the limitations of experimental high-throughput methods, computational approaches for functional module identification have emerged as a promising alternative. PPI networks are particularly well-suited for this task, as they capture not only the topology of protein interactions but also underlying functional relationships. Proteins involved in similar biological processes tend to interact more frequently, forming dense clusters within the network characterized by high intra-group connectivity and sparse inter-group interactions. Consequently, the identification of protein complexes can be formally framed as a network clustering problem, utilizing the topological properties of PPI networks to infer modular biological structures.

Beyond network topology, there is a vast repository of biological knowledge regarding proteins and the genes that encode them. Given the inherent noise and errors in experimental PPI networks, incorporating this biological information can significantly guide and refine the identification of functional modules. Consequently, computational methods are often categorized based on their ability to integrate biological data with topological features. Among the various sources of biological data, Gene Ontology (GO) is the most widely utilized, though other sources such as gene expression profiles and mutation data which also have also been successfully integrated into recent studies. GO provides a standardized, structured vocabulary for defining protein functions across three non-overlapping domains: Molecular Function (MF), Biological Process (BP), and Cellular Component (CC). By organizing these terms into a hierarchical Directed Acyclic Graph (DAG), GO enables the quantitative assessment of functional similarity between proteins, offering a robust biological context that complements the raw connectivity of PPI networks.

Although the integration of biological data significantly enhances the identification of functional modules, a substantial body of research focuses exclusively on network topology. These topology-based approaches often achieve competitive performance and are typically classified into two primary categories: divisive and agglomerative methods. Divisive methods iteratively partition the global PPI network into smaller sub-networks to isolate functional modules. Conversely, agglomerative methods often referred to as expansion or seeding approaches which initiate from seed nodes and iteratively merge or expand these sets to delineate functional modules.

Despite the rapid advancements in deep learning (DL) and graph neural networks (GNNs) across various domains, their application to functional module identification remains relatively unexplored. GNNs are particularly well-suited for this task, as they can simultaneously encode both network topology and heterogeneous protein features, effectively framing the problem as attributed graph clustering. However, the optimal integration of protein features and network structure within GNN architectures remains an open research question. Moreover, existing GNN-based approaches typically rely on two-stage pipelines, where node representations are learned first and clustering is performed as a separate post-processing step. This disconnect prevents the model from generating ‘clustering-aware’ representations that directly optimize for modular structure. To address these gaps, this study proposes YYYY, an end-to-end deep learning framework designed to integrate representation learning and community detection using GNNs into a unified optimization process, thereby enabling the discovery of biologically meaningful functional modules.

The remainder of this paper is organized as follows. Section 2 summarizes relevant studies on protein complex identification within PPI networks. Section 3 outlines the necessary preliminaries and describes the proposed methodology in detail. In Section 4, we define the evaluation metrics and describe the datasets utilized for our experiments. Section 5 presents the experimental results, providing a comprehensive comparison against existing benchmarks and state-of-the-art methods. Finally, Section 6 concludes the paper and discusses potential directions for future work.